\section{Induksi pada Pertidaksamaan}

Induksi matematika tidak terbatas pada pembuktian persamaan. Teknik ini juga sangat efektif untuk membuktikan \logic{pertidaksamaan} yang melibatkan bilangan bulat. Bukti induksi pada pertidaksamaan memerlukan kehati-hatian tambahan, karena langkah induksi harus menunjukkan bahwa ketidaksamaan ``tetap terjaga'' ketika $n$ bertambah dari $k$ ke $k+1$ \cite{rosen2019}.

\subsection{Pertidaksamaan Eksponensial}

\begin{contoh}
\textbf{Teorema}: Untuk setiap bilangan bulat $n \geq 1$, berlaku $2^n > n$.

\textbf{Bukti (Induksi):}

\textbf{Langkah Basis} ($n = 1$):
\[ 2^1 = 2 > 1 = n. \quad \checkmark \]

\textbf{Hipotesis Induksi}: Asumsikan untuk $n = k \geq 1$:
\[ 2^k > k \]

\textbf{Langkah Induksi}: Buktikan $2^{k+1} > k + 1$.
\begin{align*}
2^{k+1} &= 2 \cdot 2^k \\
         &> 2 \cdot k  & \text{(hipotesis induksi: } 2^k > k) \\
         &= k + k \\
         &\geq k + 1 & \text{(karena } k \geq 1)
\end{align*}

Sehingga $2^{k+1} > k + 1$. $\blacksquare$
\end{contoh}

\subsection{Pertidaksamaan Faktorial}

\begin{contoh}
\textbf{Teorema}: Untuk setiap bilangan bulat $n \geq 4$, berlaku $n! > 2^n$.

\textbf{Bukti (Induksi):}

\textbf{Langkah Basis} ($n = 4$):
\[ 4! = 24 > 16 = 2^4. \quad \checkmark \]

Perhatikan bahwa basis dimulai dari $n = 4$, bukan $n = 1$, karena pertidaksamaan tidak berlaku untuk $n = 1, 2, 3$ (misalnya $3! = 6 < 8 = 2^3$).

\textbf{Hipotesis Induksi}: Asumsikan untuk $n = k \geq 4$:
\[ k! > 2^k \]

\textbf{Langkah Induksi}: Buktikan $(k+1)! > 2^{k+1}$.
\begin{align*}
(k+1)! &= (k+1) \cdot k! \\
        &> (k+1) \cdot 2^k & \text{(hipotesis induksi)} \\
        &> 2 \cdot 2^k & \text{(karena } k+1 \geq 5 > 2 \text{ untuk } k \geq 4) \\
        &= 2^{k+1}
\end{align*}

Sehingga $(k+1)! > 2^{k+1}$. $\blacksquare$
\end{contoh}

\subsection{Pertidaksamaan Bernoulli}

\begin{contoh}
\textbf{Teorema (Pertidaksamaan Bernoulli)}: Untuk setiap bilangan real $x \geq -1$ dan setiap bilangan bulat $n \geq 1$:
\[ (1 + x)^n \geq 1 + nx \]

\textbf{Bukti (Induksi pada $n$):}

\textbf{Langkah Basis} ($n = 1$):
\[ (1+x)^1 = 1 + x = 1 + 1 \cdot x. \quad \checkmark \]

\textbf{Hipotesis Induksi}: Asumsikan untuk $n = k \geq 1$:
\[ (1 + x)^k \geq 1 + kx \]

\textbf{Langkah Induksi}: Buktikan $(1 + x)^{k+1} \geq 1 + (k+1)x$.
\begin{align*}
(1 + x)^{k+1} &= (1 + x)^k \cdot (1 + x) \\
               &\geq (1 + kx)(1 + x) & \text{(hip. induksi; } 1+x \geq 0) \\
               &= 1 + x + kx + kx^2 \\
               &= 1 + (k+1)x + kx^2 \\
               &\geq 1 + (k+1)x & \text{(karena } kx^2 \geq 0)
\end{align*}

Sehingga $(1+x)^{k+1} \geq 1 + (k+1)x$. $\blacksquare$
\end{contoh}

\begin{catatan}
Pada langkah induksi pertidaksamaan, perhatikan bahwa kita sering memerlukan \textbf{fakta tambahan} selain hipotesis induksi --- misalnya ``$k \geq 1$'' atau ``$kx^2 \geq 0$''. Fakta tambahan ini harus berasal dari batasan domain yang diberikan dalam pernyataan teorema, dan harus dinyatakan secara eksplisit dalam bukti.
\end{catatan}
